%% Description of the DCASE 2022 Task 2, bearing dataset.
%% Covers dimensions, combinations of faults/bearings/conditions, EDA findings,
%% supported tests and splits; the test rig is not described here.
%% Build (from eda/dcase2022_bearing/):   pdflatex dcase2022_bearing && pdflatex dcase2022_bearing

\documentclass[11pt,a4paper]{article}
\usepackage[T1]{fontenc}
\usepackage[utf8]{inputenc}
\usepackage[margin=2.4cm]{geometry}
\usepackage{amsmath}
\usepackage{booktabs}
\usepackage{graphicx}
\usepackage[hidelinks]{hyperref}

\title{DCASE 2022 Task 2, bearing}
\author{}
\date{\today}

\begin{document}
\maketitle

\section{Source}

\section{Dimensions of the time series}

\section{Combinations of faults, bearings and operating conditions}

\begin{table}[h]
\centering
\small
\begin{tabular}{lllcc}
\toprule
Section & Shift between domains & Train (normal): source / target & Test: normal / anomaly \\
\midrule
00 & rotation velocity   & 990 / 10 & 100 / 100 \\
01 & microphone location & 990 / 9  & 100 / 100 \\
02 & factory noise       & 990 / 10 & 100 / 100 \\
\bottomrule
\end{tabular}
\end{table}

\begin{itemize}
  \item Each number counts 10\,s clips: 16\,kHz mono, 160,000 samples, int16; 3,599 clips.
  \item Training holds normal clips only; anomalies appear only in test (anomaly detection).
  \item Test clips come from both domains; the split per domain is in the captures table.
  \item Per-clip attributes (e.g.\ velocity, microphone location) are in the file name and
        the attribute CSVs. No bearing IDs and no fault types.
  \item Source: \texttt{harness/data/dcase2022\_bearing.py}, D251.
\end{itemize}

\section{Findings from EDA}

\section{Tests the data support}

\section{Splits}

\end{document}
